% ID: SEL-D0009-BA-E01
% ESTADO: cotejada. Segunda pasada visual completa y borrador compilado revisados.
% Fuente: Selectividad/matematicas_aplicadas/ORD_RESER_Modelo-3-Enunciados (Mat-Ap-2021).pdf
% SHA-256: 1dc0455361ce91a1157a7b3e4a06db95bc988e8dd0c9756e9aeb74b8cbb73a81
% Página: 1; bloque A; ejercicio original 1.

\begin{enumerate}[label=\alph*)]
  \item \textbf{(1 punto)} Una frutería vende dos tipos de surtidos de frutos
  rojos, $A$ y $B$. El surtido de tipo $A$ contiene 75 g de arándanos,
  100 g de frambuesas y se vende a 2.40 euros, mientras que el de tipo $B$
  contiene 75 g de arándanos, 50 g de frambuesas y se vende a 1.80 euros.
  La frutería dispone de un total de 3.75 kg de arándanos y 4 kg de frambuesas
  y el número de surtidos que vende del tipo $A$, siempre es menor o igual
  al doble de los del tipo $B$.

  Formule, \underline{sin resolver}, el problema que permite obtener el número
  de surtidos de cada tipo que debe vender para que el beneficio sea máximo.
  \item \textbf{(1.5 puntos)} Represente el recinto limitado por las siguientes
  restricciones, calculando sus vértices:
  \[
  x+4y\geq 5\qquad x+2y\geq 4\qquad 7x+5y\leq 35\qquad x\geq 0
  \]
  \textquestiondown{}En qué punto de la región anterior la función
  $F(x,y)=2x+y$ alcanza el mínimo y cuál es dicho valor?
\end{enumerate}
